% --- Archon physics formalization source begin ---
% archon:physics
% archon:covers PhyXMiniProblems/problem_phyx_mini_0650.lean
% archon:source-report reports/phyx_mini/problem_phyx_mini_0650.source.json

\chapter{Physics problem phyx\_mini\_0650}
\label{ch:PhyXMiniProblems_problem_phyx_mini_0650}

\paragraph{Problem source.}
\#\# Physical scenario

An observer moving at a speed of \textbackslash{}( 0.995c \textbackslash{}) relative to a rod as shown in fig measures its length to be \textbackslash{}( 2.00 \textbackslash{}text\{ m\} \textbackslash{}) and sees its length to be oriented at \textbackslash{}( 30.0\textasciicircum{}\{\textbackslash{}circ\} \textbackslash{}) with respect to its direction of motion.

\#\# Question

What is the proper length of the rod?

\#\# Answer choices

- A: A: \textbackslash{}( 12.1 \textbackslash{}text\{ m\} \textbackslash{})

- B: B: \textbackslash{}( 17.4 \textbackslash{}text\{ m\} \textbackslash{})

- C: C: \textbackslash{}( 11.5 \textbackslash{}text\{ m\} \textbackslash{})

- D: D: \textbackslash{}( 13.3 \textbackslash{}text\{ m\} \textbackslash{})

\#\# Figure caption (auxiliary; use the image as primary evidence)

The image shows a rod positioned at an angle above a dashed horizontal line. The rod is labeled "2.00 m," indicating its length, and it's inclined at a "30.0°" angle with respect to the horizontal line. There is a large arrow labeled "Motion of observer" pointing to the right, below the dashed line.

\#\# Dataset metadata

Relativity; Physical Model Grounding Reasoning, Multi-Formula Reasoning

\paragraph{Recorded answer/context.}
B: B: \textbackslash{}( 17.4 \textbackslash{}text\{ m\} \textbackslash{})

\paragraph{Figure/image path.}
/root/proposal\_for\_physic/science-mango/phyx\_mini\_run/phyx\_data/test\_image/650.png

\paragraph{Formalization target.}
create a compiling Lean file with sorry bodies at `PhyXMiniProblems/problem\_phyx\_mini\_0650.lean`. The Lean declarations must preserve the physical quantities, dimensions or dimensional roles, figure labels, governing-law hypotheses, and final relation expressed by this problem.
Use Mathlib/Physlib names found through LeanExplore where available. If a physics API is missing, introduce faithful local abstractions rather than scalar placeholder aliases.

\begin{theorem}[Physics formalization target]
\label{thm:physics:phyx_mini_0650:target}
\lean{PhyXMiniProblems.ProblemPhyXMini0650.problem_phyx_mini_0650}
\uses{def:physics:phyx-mini-0650:phyxminiproblems-problemphyxmini0650-obliquerelativisticrodsetup, def:physics:phyx-mini-0650:phyxminiproblems-problemphyxmini0650-matchesobliquerodscenario, def:physics:phyx-mini-0650:phyxminiproblems-problemphyxmini0650-matchessuppliedimage650, def:physics:phyx-mini-0650:phyxminiproblems-problemphyxmini0650-matchesproblemstatementreadouts, def:physics:phyx-mini-0650:phyxminiproblems-problemphyxmini0650-hasphysicalobliquerodparameters, def:physics:phyx-mini-0650:phyxminiproblems-problemphyxmini0650-satisfiesrodcomponentgeometry, def:physics:phyx-mini-0650:phyxminiproblems-problemphyxmini0650-satisfiesrelativisticrodlengthtransformation, def:physics:phyx-mini-0650:phyxminiproblems-problemphyxmini0650-matchesdisplayedproperlengthchoice}
The assigned autoformalize agent should translate this physics problem into Lean declarations in the covered file, with theorem and lemma proof bodies written as `by sorry`.
\end{theorem}
\begin{proof}
This is an autoformalization task, not a proof task. Produce statements that can later be proved by the physics prover without weakening the physical model.
\end{proof}

% --- Archon declaration topology begin ---
\section{Declaration topology}
\begin{definition}[Length Magnitude]
  \label{def:physics:phyx-mini-0650:phyxminiproblems-problemphyxmini0650-lengthmagnitude}
  \lean{PhyXMiniProblems.ProblemPhyXMini0650.LengthMagnitude}
  A nonnegative, unit-independent physical length magnitude
\end{definition}
\begin{proof}
  This is the declared model definition of Length Magnitude.
\end{proof}
\begin{definition}[Signed Length]
  \label{def:physics:phyx-mini-0650:phyxminiproblems-problemphyxmini0650-signedlength}
  \lean{PhyXMiniProblems.ProblemPhyXMini0650.SignedLength}
  A signed, unit-independent component of a spatial displacement
\end{definition}
\begin{proof}
  This is the declared model definition of Signed Length.
\end{proof}
\begin{definition}[length Readout]
  \label{def:physics:phyx-mini-0650:phyxminiproblems-problemphyxmini0650-lengthreadout}
  \lean{PhyXMiniProblems.ProblemPhyXMini0650.lengthReadout}
  \uses{def:physics:phyx-mini-0650:phyxminiproblems-problemphyxmini0650-lengthmagnitude}
  Read a physical length magnitude in a selected length unit
\end{definition}
\begin{proof}
  This is the declared model definition of length Readout.
\end{proof}
\begin{definition}[signed Length Readout]
  \label{def:physics:phyx-mini-0650:phyxminiproblems-problemphyxmini0650-signedlengthreadout}
  \lean{PhyXMiniProblems.ProblemPhyXMini0650.signedLengthReadout}
  \uses{def:physics:phyx-mini-0650:phyxminiproblems-problemphyxmini0650-signedlength}
  Read a signed displacement component in a selected length unit
\end{definition}
\begin{proof}
  This is the declared model definition of signed Length Readout.
\end{proof}
\begin{definition}[length In Meters]
  \label{def:physics:phyx-mini-0650:phyxminiproblems-problemphyxmini0650-lengthinmeters}
  \lean{PhyXMiniProblems.ProblemPhyXMini0650.lengthInMeters}
  \uses{def:physics:phyx-mini-0650:phyxminiproblems-problemphyxmini0650-lengthmagnitude, def:physics:phyx-mini-0650:phyxminiproblems-problemphyxmini0650-lengthreadout}
  Meter readout used by the scenario and answer choices
\end{definition}
\begin{proof}
  This is the declared model definition of length In Meters.
\end{proof}
\begin{definition}[speed In Meters Per Second]
  \label{def:physics:phyx-mini-0650:phyxminiproblems-problemphyxmini0650-speedinmeterspersecond}
  \lean{PhyXMiniProblems.ProblemPhyXMini0650.speedInMetersPerSecond}
  Read a physical speed in meters per second
\end{definition}
\begin{proof}
  This is the declared model definition of speed In Meters Per Second.
\end{proof}
\begin{definition}[vacuum Speed Of Light In Meters Per Second]
  \label{def:physics:phyx-mini-0650:phyxminiproblems-problemphyxmini0650-vacuumspeedoflightinmeterspersecond}
  \lean{PhyXMiniProblems.ProblemPhyXMini0650.vacuumSpeedOfLightInMetersPerSecond}
  Physlib's exact vacuum light speed, read in meters per second
\end{definition}
\begin{proof}
  This is the declared model definition of vacuum Speed Of Light In Meters Per Second.
\end{proof}
\begin{definition}[angle Of Degrees]
  \label{def:physics:phyx-mini-0650:phyxminiproblems-problemphyxmini0650-angleofdegrees}
  \lean{PhyXMiniProblems.ProblemPhyXMini0650.angleOfDegrees}
  Convert a degree readout to Mathlib's angle type
\end{definition}
\begin{proof}
  This is the declared model definition of angle Of Degrees.
\end{proof}
\begin{definition}[Inertial Frame Label]
  \label{def:physics:phyx-mini-0650:phyxminiproblems-problemphyxmini0650-inertialframelabel}
  \lean{PhyXMiniProblems.ProblemPhyXMini0650.InertialFrameLabel}
  The rod rest frame and the inertial frame of the moving observer
\end{definition}
\begin{proof}
  This is the declared model definition of Inertial Frame Label.
\end{proof}
\begin{definition}[Axial Direction]
  \label{def:physics:phyx-mini-0650:phyxminiproblems-problemphyxmini0650-axialdirection}
  \lean{PhyXMiniProblems.ProblemPhyXMini0650.AxialDirection}
  Positive and negative directions along the relative-motion axis
\end{definition}
\begin{proof}
  This is the declared model definition of Axial Direction.
\end{proof}
\begin{definition}[Probability Plot Axis Label]
  \label{def:physics:phyx-mini-0650:phyxminiproblems-problemphyxmini0650-probabilityplotaxislabel}
  \lean{PhyXMiniProblems.ProblemPhyXMini0650.ProbabilityPlotAxisLabel}
  Literal axis labels visible in the supplied (mismatched) raster
\end{definition}
\begin{proof}
  This is the declared model definition of Probability Plot Axis Label.
\end{proof}
\begin{definition}[Inverse Length Unit Label]
  \label{def:physics:phyx-mini-0650:phyxminiproblems-problemphyxmini0650-inverselengthunitlabel}
  \lean{PhyXMiniProblems.ProblemPhyXMini0650.InverseLengthUnitLabel}
  Unit label printed on the vertical probability-density axis
\end{definition}
\begin{proof}
  This is the declared model definition of Inverse Length Unit Label.
\end{proof}
\begin{definition}[Probability Plot Feature]
  \label{def:physics:phyx-mini-0650:phyxminiproblems-problemphyxmini0650-probabilityplotfeature}
  \lean{PhyXMiniProblems.ProblemPhyXMini0650.ProbabilityPlotFeature}
  Named graphical features visible in image 650
\end{definition}
\begin{proof}
  This is the declared model definition of Probability Plot Feature.
\end{proof}
\begin{definition}[Supplied Probability Density Figure]
  \label{def:physics:phyx-mini-0650:phyxminiproblems-problemphyxmini0650-suppliedprobabilitydensityfigure}
  \lean{PhyXMiniProblems.ProblemPhyXMini0650.SuppliedProbabilityDensityFigure}
  \uses{def:physics:phyx-mini-0650:phyxminiproblems-problemphyxmini0650-probabilityplotaxislabel, def:physics:phyx-mini-0650:phyxminiproblems-problemphyxmini0650-inverselengthunitlabel, def:physics:phyx-mini-0650:phyxminiproblems-problemphyxmini0650-probabilityplotfeature}
  Typed scalar readout of the actual image. The function is explicitly a plotted value in inverse centimeters, rather than a replacement for a dimensionful rod quantity
\end{definition}
\begin{proof}
  This is the declared model definition of Supplied Probability Density Figure.
\end{proof}
\begin{definition}[Oblique Relativistic Rod Setup]
  \label{def:physics:phyx-mini-0650:phyxminiproblems-problemphyxmini0650-obliquerelativisticrodsetup}
  \lean{PhyXMiniProblems.ProblemPhyXMini0650.ObliqueRelativisticRodSetup}
  \uses{def:physics:phyx-mini-0650:phyxminiproblems-problemphyxmini0650-lengthmagnitude, def:physics:phyx-mini-0650:phyxminiproblems-problemphyxmini0650-signedlength, def:physics:phyx-mini-0650:phyxminiproblems-problemphyxmini0650-inertialframelabel, def:physics:phyx-mini-0650:phyxminiproblems-problemphyxmini0650-axialdirection, def:physics:phyx-mini-0650:phyxminiproblems-problemphyxmini0650-suppliedprobabilitydensityfigure}
  The independent physical observables used in the rod problem. The proper length and all rest-frame components are fields, not definitions made from an answer choice or from the requested numerical result
\end{definition}
\begin{proof}
  This is the declared model definition of Oblique Relativistic Rod Setup.
\end{proof}
\begin{definition}[speed Fraction Of Light]
  \label{def:physics:phyx-mini-0650:phyxminiproblems-problemphyxmini0650-speedfractionoflight}
  \lean{PhyXMiniProblems.ProblemPhyXMini0650.speedFractionOfLight}
  \uses{def:physics:phyx-mini-0650:phyxminiproblems-problemphyxmini0650-speedinmeterspersecond, def:physics:phyx-mini-0650:phyxminiproblems-problemphyxmini0650-vacuumspeedoflightinmeterspersecond, def:physics:phyx-mini-0650:phyxminiproblems-problemphyxmini0650-obliquerelativisticrodsetup}
  The dimensionless relative speed β = v/c
\end{definition}
\begin{proof}
  This is the declared model definition of speed Fraction Of Light.
\end{proof}
\begin{definition}[lorentz Factor]
  \label{def:physics:phyx-mini-0650:phyxminiproblems-problemphyxmini0650-lorentzfactor}
  \lean{PhyXMiniProblems.ProblemPhyXMini0650.lorentzFactor}
  \uses{def:physics:phyx-mini-0650:phyxminiproblems-problemphyxmini0650-obliquerelativisticrodsetup, def:physics:phyx-mini-0650:phyxminiproblems-problemphyxmini0650-speedfractionoflight}
  Physlib's Lorentz factor for the setup's dimensionless speed
\end{definition}
\begin{proof}
  This is the declared model definition of lorentz Factor.
\end{proof}
\begin{definition}[Matches Oblique Rod Scenario]
  \label{def:physics:phyx-mini-0650:phyxminiproblems-problemphyxmini0650-matchesobliquerodscenario}
  \lean{PhyXMiniProblems.ProblemPhyXMini0650.MatchesObliqueRodScenario}
  \uses{def:physics:phyx-mini-0650:phyxminiproblems-problemphyxmini0650-obliquerelativisticrodsetup}
  Frame and direction assignments stated by the rod scenario
\end{definition}
\begin{proof}
  This is the declared model definition of Matches Oblique Rod Scenario.
\end{proof}
\begin{definition}[Matches Supplied Image650]
  \label{def:physics:phyx-mini-0650:phyxminiproblems-problemphyxmini0650-matchessuppliedimage650}
  \lean{PhyXMiniProblems.ProblemPhyXMini0650.MatchesSuppliedImage650}
  \uses{def:physics:phyx-mini-0650:phyxminiproblems-problemphyxmini0650-suppliedprobabilitydensityfigure}
  Literal evidence in the actual primary image: x is in centimeters, |ψ(x)|² is in inverse centimeters, and the green graph is |x| on [-1,1] with vertical endpoint drops and zero exterior. The final two fields make the source mismatch explicit
\end{definition}
\begin{proof}
  This is the declared model definition of Matches Supplied Image650.
\end{proof}
\begin{definition}[Matches Problem Statement Readouts]
  \label{def:physics:phyx-mini-0650:phyxminiproblems-problemphyxmini0650-matchesproblemstatementreadouts}
  \lean{PhyXMiniProblems.ProblemPhyXMini0650.MatchesProblemStatementReadouts}
  \uses{def:physics:phyx-mini-0650:phyxminiproblems-problemphyxmini0650-lengthinmeters, def:physics:phyx-mini-0650:phyxminiproblems-problemphyxmini0650-angleofdegrees, def:physics:phyx-mini-0650:phyxminiproblems-problemphyxmini0650-obliquerelativisticrodsetup, def:physics:phyx-mini-0650:phyxminiproblems-problemphyxmini0650-speedfractionoflight}
  Numerical readouts from the problem prose. 199/200 = 0.995; the angle and length equalities describe only what the moving observer reports
\end{definition}
\begin{proof}
  This is the declared model definition of Matches Problem Statement Readouts.
\end{proof}
\begin{definition}[Has Physical Oblique Rod Parameters]
  \label{def:physics:phyx-mini-0650:phyxminiproblems-problemphyxmini0650-hasphysicalobliquerodparameters}
  \lean{PhyXMiniProblems.ProblemPhyXMini0650.HasPhysicalObliqueRodParameters}
  \uses{def:physics:phyx-mini-0650:phyxminiproblems-problemphyxmini0650-lengthreadout, def:physics:phyx-mini-0650:phyxminiproblems-problemphyxmini0650-signedlengthreadout, def:physics:phyx-mini-0650:phyxminiproblems-problemphyxmini0650-obliquerelativisticrodsetup, def:physics:phyx-mini-0650:phyxminiproblems-problemphyxmini0650-speedfractionoflight}
  Positivity and the subluminal domain required by the physical model
\end{definition}
\begin{proof}
  This is the declared model definition of Has Physical Oblique Rod Parameters.
\end{proof}
\begin{definition}[Satisfies Rod Component Geometry]
  \label{def:physics:phyx-mini-0650:phyxminiproblems-problemphyxmini0650-satisfiesrodcomponentgeometry}
  \lean{PhyXMiniProblems.ProblemPhyXMini0650.SatisfiesRodComponentGeometry}
  \uses{def:physics:phyx-mini-0650:phyxminiproblems-problemphyxmini0650-lengthreadout, def:physics:phyx-mini-0650:phyxminiproblems-problemphyxmini0650-signedlengthreadout, def:physics:phyx-mini-0650:phyxminiproblems-problemphyxmini0650-obliquerelativisticrodsetup}
  Euclidean component geometry for the independently measured length vectors. The observer-frame orientation determines its two components, while the proper magnitude is the norm of the independent rest-frame components
\end{definition}
\begin{proof}
  This is the declared model definition of Satisfies Rod Component Geometry.
\end{proof}
\begin{definition}[Satisfies Relativistic Rod Length Transformation]
  \label{def:physics:phyx-mini-0650:phyxminiproblems-problemphyxmini0650-satisfiesrelativisticrodlengthtransformation}
  \lean{PhyXMiniProblems.ProblemPhyXMini0650.SatisfiesRelativisticRodLengthTransformation}
  \uses{def:physics:phyx-mini-0650:phyxminiproblems-problemphyxmini0650-signedlengthreadout, def:physics:phyx-mini-0650:phyxminiproblems-problemphyxmini0650-obliquerelativisticrodsetup, def:physics:phyx-mini-0650:phyxminiproblems-problemphyxmini0650-lorentzfactor}
  The generic special-relativistic component law: only the component parallel to the boost is divided by γ; the perpendicular component is unchanged. It contains no numerical proper length and no answer-choice value
\end{definition}
\begin{proof}
  This is the declared model definition of Satisfies Relativistic Rod Length Transformation.
\end{proof}
\begin{definition}[Answer Choice]
  \label{def:physics:phyx-mini-0650:phyxminiproblems-problemphyxmini0650-answerchoice}
  \lean{PhyXMiniProblems.ProblemPhyXMini0650.AnswerChoice}
  Labels of the four source answer choices
\end{definition}
\begin{proof}
  This is the declared model definition of Answer Choice.
\end{proof}
\begin{definition}[displayed Proper Length Meters]
  \label{def:physics:phyx-mini-0650:phyxminiproblems-problemphyxmini0650-displayedproperlengthmeters}
  \lean{PhyXMiniProblems.ProblemPhyXMini0650.displayedProperLengthMeters}
  \uses{def:physics:phyx-mini-0650:phyxminiproblems-problemphyxmini0650-answerchoice}
  Displayed proper-length choices in meters, represented exactly
\end{definition}
\begin{proof}
  This is the declared model definition of displayed Proper Length Meters.
\end{proof}
\begin{definition}[recorded Dataset Answer]
  \label{def:physics:phyx-mini-0650:phyxminiproblems-problemphyxmini0650-recordeddatasetanswer}
  \lean{PhyXMiniProblems.ProblemPhyXMini0650.recordedDatasetAnswer}
  \uses{def:physics:phyx-mini-0650:phyxminiproblems-problemphyxmini0650-answerchoice}
  The answer label recorded by the source dataset
\end{definition}
\begin{proof}
  This is the declared model definition of recorded Dataset Answer.
\end{proof}
\begin{definition}[Matches Displayed Proper Length Choice]
  \label{def:physics:phyx-mini-0650:phyxminiproblems-problemphyxmini0650-matchesdisplayedproperlengthchoice}
  \lean{PhyXMiniProblems.ProblemPhyXMini0650.MatchesDisplayedProperLengthChoice}
  \uses{def:physics:phyx-mini-0650:phyxminiproblems-problemphyxmini0650-lengthinmeters, def:physics:phyx-mini-0650:phyxminiproblems-problemphyxmini0650-obliquerelativisticrodsetup, def:physics:phyx-mini-0650:phyxminiproblems-problemphyxmini0650-answerchoice, def:physics:phyx-mini-0650:phyxminiproblems-problemphyxmini0650-displayedproperlengthmeters}
  Agreement with a proper length displayed to the nearest tenth of a meter. The half-tenth tolerance is 0.05 m = 1/20 m
\end{definition}
\begin{proof}
  This is the declared model definition of Matches Displayed Proper Length Choice.
\end{proof}
% --- Archon declaration topology end ---
% --- Archon physics formalization source end ---
